\clearpage\thispagestyle{empty}
\pdfbookmark[0]{Quellen zu Mitteilung über Goldschmidts Kometenelemente (17. August 1844)}{cometaug1844notice-sources}
\begin{center}\textbf{Quellen zu Mitteilung über Goldschmidts Kometenelemente (17. August 1844)}\end{center}
\noindent Diese KI-Übertragung ins heutige Deutsch folgt der eingeklammerten Mitteilung auf S.~467 der \href{https://gdz.sub.uni-goettingen.de/id/PPN236007467}{\textit{Werke}, Band~VI} (Göttingen, 1874). Sie verweist auf den Brief, den Carl Friedrich Gauß am 28. Juli 1844 an Heinrich Christian Schumacher schrieb. Der Brief erschien in den \textit{Astronomischen Nachrichten}, Band~XXII, Nr.~512, \href{https://api.digitale-sammlungen.de/iiif/image/v2/bsb10049447_00066/full/1800,/0/default.jpg}{Sp.~115--116}. Die Quellenangabe der Werkausgabe datiert das Heft auf den 17. August 1844.
\par\smallskip\noindent Der Brief nennt Dr. Goldschmidt als Urheber der Berechnung der parabolischen Elemente. Im Erstdruck stehen die Elemente, die Abweichungen von drei Beobachtungen, die Ephemeride und die Koordinatenausdrücke. Die Werkausgabe enthält die eingeklammerte Mitteilung. KI-Übertragung und KI-Satz.
\clearpage
